% ================================================================== 2
\newpage
\section{Introducción}
\label{sec:introduccion}

El tema asignado al grupo es la \textbf{biblioteca educativa de teoría de
números y cifrado}, con el título \emph{«Optimización del pipeline de ejecución
para el sistema criptográfico RSA»}. La biblioteca está dirigida a estudiantes
que necesitan comprender los fundamentos matemáticos de RSA: en lugar de usarlo
como caja negra, permite ejecutar cada algoritmo por separado, observar los
valores intermedios y comparar el costo de distintas alternativas.

RSA ejecuta un \emph{pipeline} de etapas: generar primos grandes, construir las
claves, cifrar y descifrar. El problema es \textbf{determinar cuánto se reduce
el costo de ese pipeline al reemplazar en cada etapa una alternativa base por
una optimizada, sin alterar el resultado}. El objetivo es implementar y
comparar experimentalmente ambas alternativas en cada etapa, verificar su
corrección y contrastar los resultados con el análisis de complejidad.

El alcance comprende la aritmética modular, las pruebas de primalidad, la
factorización entera como ataque y el RSA educativo con claves según
FIPS~186-5~\cite{fips186}. No incluye el relleno OAEP~\cite{rfc8017} ni el
uso en producción.

El documento presenta los fundamentos (Sección~\ref{sec:fundamentos}), el
stack tecnológico (\ref{sec:stack}), la formulación (\ref{sec:formulacion}), el
diseño algorítmico (\ref{sec:diseno-algoritmico}), el sistema
(\ref{sec:sistema}), las pruebas (\ref{sec:pruebas}), la metodología
experimental (\ref{sec:metodologia}), los resultados (\ref{sec:resultados}), la
discusión (\ref{sec:discusion}), las conclusiones (\ref{sec:conclusiones}) y
los anexos (\ref{sec:anexos}).
